\documentclass{beamer}

\usepackage{listings}
\usepackage{xcolor}
\usepackage{amsmath, amssymb, amsfonts}

\DeclareMathOperator{\range}{Range}
\DeclareMathOperator{\nul}{Null}

\lstset{
    language=C++,
    basicstyle=\ttfamily\small,
    keywordstyle=\color{blue}\bfseries,
    commentstyle=\color{teal}
}

\usetheme{metropolis}

\title{Multigrid methods}
\subtitle{a short introduction}
\author{Gabriel Clave}
\date{\today}

\begin{document}

\frame{\titlepage}

% slide Notations
\begin{frame}{Notations}
\subsection*{Notations}
Let $A \in \mathbb{R}^{n \times n}$ be non-singular and $f \in \mathbb{R}^n$.
\begin{align*}
    \text{Problem:}          \quad & Au = f \\
    \text{Error:}            \quad & e = u - v \\
    \text{Residual:}         \quad & r = f - Av \\
    \text{Residual Problem:} \quad & Ae = r
\end{align*}
\end{frame}

\section{Relaxation methods}

% slide the smoother
% at the origin of multigrid are the relaxation methods
\begin{frame}{Relaxation Methods}
\subsection*{Relaxation Methods}
Decompose $A = M - N$, where $M$ is easy to invert:
\[
    Au = f \implies Mu = Nu + f
\]

$u$ is a fixed point of the operator:
\[
    u = M^{-1}Nu + M^{-1}f
\]
% this gives us an idea for an iteration
With iteration matrix $S = M^{-1}N = I - M^{-1}A$:
\[
    v^{(k+1)} = Sv^{(k)} + M^{-1}f
\]

Using $M^{-1}f = u - Su$, the error evolution satisfies:
\[
    e^{(k+1)} = Se^{(k)} = S^{k}e^{(0)}
\]
If $\rho(S) < 1$, the error is forced to zero.
\end{frame}

% slide two classic examples
\begin{frame}{Classic Relaxation Examples}
\subsubsection*{Examples (for matrices with nonzero diagonal)}
\begin{itemize}
    \item \textbf{Jacobi:} Split $A = D - L - U$ with $M = D$.
    \[
        Dv^{(k+1)} = (L + U)v^{(k)} + f
    \]
    \[
        a_{ii}v^{(k+1)}_{i} = \sum_{j \neq i} -a_{ij}v^{(k)}_{j} + f_{i}
    \]
    Solve for $v_{i}^{(k+1)}$ using only the previous estimate.

    \vspace{0.5em}
    \item \textbf{Gauss-Seidel:} Choose $M = D - L$.
    \[
        (D-L)v^{(k+1)} = Uv^{(k)} + f
    \]
    Solve for $v_{i}^{(k+1)}$ using already computed values at step $k+1$.
\end{itemize}
\end{frame}

% slide Smoother alone aren't great
\begin{frame}{Issues and Convergence Rates}
\subsubsection*{Issue \& Strength}
\[
    e^{(k+1)} = Se^{(k)} = S^{k}e^{(0)}
\]

\begin{itemize}
    \item Error components along eigenvectors corresponding to large eigenvalues of $S$ decay very slowly.
    \item Convergence rate stagnates quickly after initial iterations.
\end{itemize}

\vspace{1em}
\begin{center}
    % [Insert picture: Convergence stagnation / Error norm vs iteration]
\end{center}
\end{frame}
% used alone, relaxation methods are slow and expensive 

% slide their strength lies in their ability to remove "high frequency error"
\begin{frame}{Smoothing Property}
\begin{itemize}
    \item Relaxation methods are effective at damping \textbf{high-frequency / oscillatory} components of the error.
    \item They are inefficient at eliminating \textbf{low-frequency / smooth} error components.
\end{itemize}
% for classic smoother, people realized it was the "oscillating" error that was effectively removed

\vspace{2em}
\begin{center}
    % [Insert picture: High-frequency vs Low-frequency error smoothing]
\end{center}
\end{frame}

% slide coarse grid to artificially increase the frequency
\begin{frame}{Coarse Grid Concept}
\subsubsection*{Coarse Grid}

\textbf{Core Intuition:} Subsampling a smooth signal on a coarser grid turns low-frequency modes into high-frequency modes relative to the new mesh size.

\vspace{2em}
\begin{center}
    % [Insert picture: Effect of subsampling on mode frequencies]
\end{center}
\end{frame}

% slide idea
\begin{frame}{Multigrid Idea \& Grid Transfers}
\textbf{Recursive Procedure:}
\begin{enumerate}
    \item Apply relaxation on the fine grid until the error is smooth.
    \item Project the smooth residual/error onto a coarser grid where it becomes oscillatory.
    \item Perform relaxation on the coarse grid.
    \item Correct the fine grid approximation.
\end{enumerate}

\vspace{1em}
\textbf{Inter-grid Transfer Operators:}
\begin{itemize}
    \item \textbf{Interpolation / Prolongation ($P$):} Coarse $\to$ Fine
    \item \textbf{Projection / Restriction ($R = P^\top$):} Fine $\to$ Coarse
\end{itemize}
\end{frame}

% slide coarse grid
\begin{frame}{Coarse Grid Operator Formulation}
\subsubsection*{Coarse Grid Equations}
On the coarse grid, we solve the residual equation:
\[
    A_{c}e_{c} = r_{c}
\]

\begin{itemize}
    \item \textbf{Geometric Case:} The underlying continuous physics dictate the operator; $A_c$ is discretized directly on the coarse geometry.
    \item \textbf{Algebraic Case (Galerkin approach):} Given prolongation $P$ and setting $e \approx Pe_{c}$:
\end{itemize}
\begin{align*}
    Ae = r &\implies APe_{c} \approx r \\ 
           &\implies P^{\top}APe_{c} = P^{\top}r
\end{align*}

Choosing $r_c = P^\top r$ yields the Galerkin coarse-grid operator:
\[
    A_{c} = P^{\top}AP
\]
\end{frame}

\section{Basic Multigrid}

% slide general idea
\begin{frame}{Two-Grid Correction Scheme}
\subsection*{General Idea}

\begin{itemize}
    \item \textbf{Pre-smoothing:} Relax on $Au = f$ on the fine grid to obtain an approximation $v$.
    \item \textbf{Residual Computation:} Compute $r = f - Av$.
    \item \textbf{Coarse Grid Solve:}
    \begin{itemize}
        \item Solve on the residual equation $A_{c}e_{c} = r_{c}$ on the coarse grid to obtain an error approximation $e_{c}$.
    \end{itemize}
    \item \textbf{Coarse Grid Correction:} Correct the approximation with $v \leftarrow v + Pe_{c}$.
\end{itemize}
\end{frame}

% slide V-cycle
\begin{frame}{Multigrid V-Cycle}
\subsection*{V-Cycle}

\begin{itemize}
    \item \textbf{Pre-smoothing:} Relax on $Au = f$ on the fine grid to obtain an approximation $v$.
    \item \textbf{Residual Computation:} Compute $r = f - Av$.
    \item \textbf{Recursion:}
    \begin{itemize}
        \item Recursively solve $A_{c}e_{c} = r_{c}$ on the coarse grid
        \item apply the same procedure down to the coarsest level.
    \end{itemize}
    \item \textbf{Coarse Grid Correction:} Correct the approximation on the fine grid: $v \leftarrow v + Pe_{c}$.
    \item \textbf{Post-smoothing:} (Optional) Apply additional relaxation steps on $Au = f$.
\end{itemize}
\end{frame}

% slide multigrid operator
\begin{frame}{Two-Grid Operator Derivation}
\subsection*{Two-Grid Operator}

\begin{align*}
    \text{Pre-smoothing ($q$ steps):} & \quad e \leftarrow S^{q}e \\
    \text{Compute residual:}          & \quad r = Ae \\
    \text{Restrict residual:}         & \quad r_{c} = P^{\top}r \\
    \text{Exact coarse solve:}        & \quad e_{c} = A_{c}^{-1}r_{c} = {(P^{\top}AP)}^{-1}r_{c} \\
    \text{Interpolate \& Correct:}    & \quad e \leftarrow e - Pe_{c}
\end{align*}

\vspace{0.5em}
Combining these steps we get:
\[
    e \leftarrow (I - P{(P^{\top}AP)}^{-1}P^{\top}A)S^{q}e
\]

The \textbf{Two-Grid Error Operator} is:
\[
    TG_{q} = (I - P{(P^{\top}AP)}^{-1}P^{\top}A)S^{q}
\]
\end{frame}

% slide algebraic insight
\begin{frame}{Algebraic Insight}
\subsection*{Algebraic Insight}

\[
    TG_{q} = (I - P{(P^{\top}AP)}^{-1}P^{\top}A)S^{q}
\]

$\Pi = P{(P^{\top}AP)}^{-1}P^{\top}A$ is the $A$-orthogonal projection onto $\range(P)$.

\begin{itemize}
    \item $TG_{0} = I - \Pi$ is the $A$-orthogonal projection onto $\range(P)^{\perp_A} = \nul(P^{\top}A)$.
    \item Direct sum decomposition: $V = \range(P) \oplus_{A} \nul(P^{\top}A)$.
    \item The coarse grid correction eliminates the error component along $\range(P)$.
\end{itemize}

\end{frame}

\begin{frame}{Algebraic Insight}

The coarse grid correction $TG_{0} = I - \Pi$ eliminates the error component along $\range(P)$.

\vspace{0.5em}
\textbf{Complementary Mechanism:}
\begin{itemize}
    \item Decomposition into Fourier modes: $V = L \oplus H$ (low and high frequencies).
    \item Relaxation $S^q$ eliminates high-frequency components ($H$).
    \item Efficiency relies on the near-alignment of low frequencies and range of prolongation:
    \[
        L \approx \range(P)
    \]
\end{itemize}
\end{frame}

% slide Spectral insight
\begin{frame}{Spectral Insight}
\subsection*{Spectral Insight}

\textbf{Action of the Two-Grid Operator on Modes:}

\begin{itemize}
    \item \textbf{Coarse Grid Correction alone ($TG_0$):}
    \begin{itemize}
        \item Suppresses smooth / low-frequency components of the error.
        \item Leaves high-frequency components largely untouched.
    \end{itemize}
    
    \vspace{0.5em}
    % picture

\end{itemize}
\end{frame}

\begin{frame}{Spectral Insight}
\subsection*{Spectral Insight}

\textbf{Action of the Two-Grid Operator on Modes:}

\begin{itemize}
    
    \item \textbf{Combined Two-Grid Operator ($TG_q = TG_0 S^q$):}
    \begin{itemize}
        \item $S^q$ damps high frequencies.
        \item $TG_0$ eliminates low frequencies.
        \item Result: Uniform suppression of all error components across the spectrum.
    \end{itemize}
\end{itemize}
\end{frame}

% slide final graph
\begin{frame}{Summary Scheme}
\subsection*{Summary Scheme}

\begin{center}
    % [Insert graph: Spectrum / Error reduction comparison]
    % e.g., Plots showing error spectrum for S^q, TG_0, and TG_q
\end{center}
\end{frame}

% \section{algebraic multi grid}

% There is no grid: location are unknown or inexistant.

% \subsection*{algebraic smooth error}

% we choose a relaxation method. An error is algebraically smooth when it is not reduced by the smoother.

% \begin{align*}
%       & Se \approx  e \\ 
%  \iff & \left(I - M^{-1}N\right)e \approx  0 \\ 
%  \iff & \left(M - N\right)e \approx  0 \\ 
%  \iff & Ae = r \approx  0 \\ 
% \end{align*}

% more rigorously if the smoother satisfies the smoothing property:

% we have $\|r\|_{D^{-1}} \ll  \|e\|_{A}$

% \begin{align*}
%     \|e\|_{A}^{2} & = \langle Ae, e \rangle \\
%                   & =  \langle D^{-\frac{1}{2}}Ae, D^{\frac{1}{2}}e \rangle \\
%                   & \le \| D^{-\frac{1}{2}}Ae \|_{2} \| D^{\frac{1}{2}}e \| \\
%                   & \le \| r \|_{D^{-1}} \| e \|_{D}
% \end{align*}

% meaning we have $ \|e\|_{A} \ll \|e\|_{D}$ or $\langle Ae,e \rangle \ll \langle De, e \rangle  $

% for some smooth error $s = D^{\frac{1}{2}}e$

% we have $ \langle Ae,e \rangle = s^{\top}D^{-\frac{1}{2}}AD^{-\frac{1}{2}}s $ and $\langle De, e \rangle = s^{\top}s $

% meaning
% \begin{align*}
%       & s^{\top}\hat{A} s \ll  s^{\top}s \\ 
%  \iff & \frac{s^{\top}\hat{A} s}{s^{\top}s} \ll 1 \\ 
%  \iff & r_{\hat{A}}(s) \ll 1 \\ 
% \end{align*}

% with $\hat{A} = D^{-\frac{1}{2}}AD^{-\frac{1}{2}}$

% the Rayleigh quotient of a smooth error is small.

% This means that smooth errors are a combination of eigenvectors of $\hat{A}$ with small eigenvalues.

% % but it is \hat{A} and not A and s and not e ?
% % TODO look at: generalized eigenvalue problem

% \textbf{Connection with the smoother}

% \[
% s^{(k+1)} = D^{1/2}\left(I - M^{-1}A\right)D^{-1/2}s^{(k)} = \left(I - \hat{M}^{-1}\hat{A}\right)s^{(k)} = \hat{S}s^{(k)}
% \]

% Damped Jacobi: $M = \frac{1}{\omega}D$
% In this case, $\hat{M} = \frac{1}{\omega}I$

% The smoother operator simplifies to: $\hat{S} = I - \omega \hat{A}$

% The eigenvalues are $\lambda_i(\hat{S}) = 1 - \omega \lambda_i(\hat{A})$.

% For large eigenvalues of $\hat{A}$, $|\lambda_i(\hat{S})| \ll 1$: the smoother is effective

% For small eigenvalues, we get $\lambda_i(\hat{S}) \approx 1$: these components form the smooth error, unaffected by the iteration

% \subsubsection*{Spectral insight}

% we want to compare $\|r\|_{D^{-1}} $ and $\|e\|_{A} = \langle Ae,e \rangle$ when $e$ is a smooth error.

% \[
% \langle Ae,e \rangle = \langle \hat{A}s,s \rangle = \sum_{k}{\lambda_{k}c_{k}^{2}}
% \]
% using an eigenvector basis of $\hat{A}$

% and
% \[
% \|r\|_{D^{-1}} = \|Ae\|_{D^{-1}} = \|D^{-\frac{1}{2}}Ae\|_{2} = \|\hat{A}s\|_{2}
% \]

% if the smoother is effective, only components of eigenvectors of $\hat{A}$ with small eigenvalues remain, ie:

% \[
% \|r\|_{D^{-1}}^{2} \approx \sum_{\lambda_{k} \le \epsilon}{\lambda_{k}^{2}c_{k}^{2}} \le \epsilon \sum_{\lambda_{k} \le \epsilon}{\lambda_{k}c_{k}^{2}} \approx \epsilon \langle Ae,e \rangle
% \]

% ie

% \[
% \|r\|_{D^{-1}} \le \sqrt{\epsilon}\|e\|_{A}
% \]

% if the smoother is effective $\sqrt{\epsilon} \ll 1$ and we have $\|r\|_{D^{-1}} \ll  \|e\|_{A}$

% \subsection*{classic AMG}


% $e_{i}$ is approximated by a weighted average of its neighbors. We focus on the biggest contribution, ie the highest $|a_{ij}|$

% strong dependence:
% \[
% |a_{ij}| \ge \theta \max_{i \ne k}{|a_{ik}|}
% \]

% if j strongly influence i, it's a good candidate for the interpolation.
% % look for a better geometric explanation

% just like in the geometric case, when the error is smooth, local (= strong influence) variation are small.

% In classic AMG, the interpolation is given by:

% \begin{equation*}
%     {(P \mathbf{e})}i = 
%     \begin{cases}
%         \hfill e_i & \text{if } i \in C, \\[8pt]
%         \displaystyle \sum_{j \in C_i} \omega_{ij} e_j & \text{if } i \in F
%     \end{cases}
% \end{equation*}

% Coarse grid values are directly used. The other are inferred from their strong connexion with the coarse grid.

% Once the interpolation is build, the restriction and coarse operator are defined as:

% \begin{align*}
%       & R = P^{\top} \\ 
%       & A_{c} = P^{\top}AP
% \end{align*}

% This way the correction is optimal in the A-norm:

% \begin{align*}
%       \|TGe\|_{A} = \|e - \Pi_{A}e\|_{A} & = \min_{w \in \range(P)}{\|e - w\|_{A}} \\ 
%                     \|e - Pe_{c}\|_{A}   & = \min_{x \in \Omega_{c}}{\|e - Px\|_{A}}
% \end{align*}


% by property of the $A$-orthogonal projection on $\range(P)$ $\Pi_{A}$

\end{document}